\setcounter{section}{2}

  \zwischenueberschrift{Soal latihan}

\inputaufgabegibtloesung
{}
{

Misalkan
\maabbdisp {f,g} {\R} { \R^n
} {}
\definitionsverweis {kurva-kurva terdiferensialkan}{}{.}
Hitung
\definitionsverweis {turunan}{}{}
fungsi
\maabbeledisp {} {\R} {\R
} { t } { \left\langle f(t) , g(t)  \right\rangle
} {.}
Nyatakan hasilnya menggunakan hasil kali dalam.

}
{} {}

\inputaufgabegibtloesung
{}
{

Misalkan
\maabb {h} { I } { V
} {}
sebuah kurva yang terdiferensialkan dua kali dalam
\definitionsverweis {ruang vektor Euklides}{}{}
$V$. Buktikan bahwa jika

\mavergleichskette
{\vergleichskette
{h'(t)
}
{ \neq }{ 0
}
{  }{
}
{    }{
}
{   }{
}
}
{}{}{,}
maka berlaku persamaan

\mavergleichskettedisp
{\vergleichskette
{ \Vert { h'(t)} \Vert'
}
{ =} { { \frac{   \left\langle h'(t) , h^{\prime \prime} (t)  \right\rangle  }{  \left\langle h'(t) , h'(t)  \right\rangle  } } \Vert { h'(t)} \Vert
}
{ } {
}
{ } {
}
{ } {
}
}
{}{}{.}

}
{} {}

Sebuah
\definitionsverweis {kurva terdiferensialkan}{}{}
\maabbeledisp {f} { [a,b] } {\R^n
} { t } {f(t) = \left( f_1(t)  , \,   \ldots  , \,  f_n(t)                 \right)
} {,}
disebut
\definitionswort {berparametrisasi panjang busur}{,}
jika

\mavergleichskette
{\vergleichskette
{ f_1'(t)^2 + \cdots +  f_n'(t)^2
}
{ = }{ 1
}
{  }{
}
{    }{
}
{   }{
}
}
{}{}{}
untuk setiap $t$.

\inputaufgabe
{}
{

Misalkan
\maabb {\gamma} { I } {\R^n
} {}
sebuah kurva yang dua kali
\definitionsverweis {terdiferensialkan secara kontinu}{}{}
dan
\definitionsverweis {berparametrisasi panjang busur}{}{.}
Buktikan bahwa
\mathkor {} {\gamma'(t)} {dan} {\gamma^{\prime \prime}(t)} {}
saling tegak lurus.

}
{} {}

\inputaufgabe
{}
{

Misalkan

\mavergleichskette
{\vergleichskette
{ P
}
{ \in }{ \R^n
}
{  }{
}
{    }{
}
{   }{
}
}
{}{}{}
sebuah titik dan misalkan

\mavergleichskette
{\vergleichskette
{ I
}
{ = }{ {]{-1},1[}
}
{  }{
}
{    }{
}
{   }{
}
}
{}{}{.}
Kita tinjau himpunan

\mavergleichskettedisp
{\vergleichskette
{ M
}
{ =} { \left\{ f:I \rightarrow \R^n   \mid  f \text{ terdiferensialkan} , \,  f(0) = P       \right\}
}
{ } {
}
{ } {
}
{ } {
}
}
{}{}{.}
Dua
\definitionsverweis {kurva}{}{}

\mavergleichskette
{\vergleichskette
{ f,g
}
{ \in }{ M
}
{  }{
}
{    }{
}
{   }{
}
}
{}{}{}
disebut \stichwort {ekuivalen secara tangensial} {,} jika

\mavergleichskettedisp
{\vergleichskette
{ f'(0)
}
{ =} { g'(0)
}
{ } {
}
{ } {
}
{ } {
}
}
{}{}{.}

a) Buktikan bahwa ini merupakan suatu
\definitionsverweis {relasi ekuivalensi}{}{.}

b) Temukan wakil paling sederhana bagi setiap
\definitionsverweis {kelas ekuivalensi}{}{.}

c) Berikan satu wakil lain untuk setiap kelas.

d) Deskripsikan himpunan kelas-kelas ekuivalensi tersebut
\zusatzklammer {yakni
\definitionsverweis {himpunan hasil bagi}{}{}} {} {.}

}
{} {}

\inputaufgabe
{}
{

Dengan syarat
\zusatzklammer {reguler di setiap titik} {} {,}
berikan sebuah contoh
\definitionsverweis {hipermuka terdiferensialkan}{}{,}
yang tidak
\definitionsverweis {terhubung}{}{.}

}
{} {}

\inputaufgabe
{}
{

Misalkan

\mavergleichskette
{\vergleichskette
{ V
}
{ \subseteq }{ \R^{n-1}
}
{  }{
}
{    }{
}
{   }{
}
}
{}{}{}
sebuah
\definitionsverweis {himpunan terbuka}{}{}
dan misalkan
\maabbdisp {f} { V } {\R
} {}
sebuah fungsi yang
\definitionsverweis {terdiferensialkan secara kontinu}{}{.}
Kita tinjau
\definitionsverweis {grafik}{}{}
$Y$ dari $f$ sebagai hipermuka terdiferensialkan di $\R^n$ dalam pengertian
Contoh 1.2.
Tentukan
\definitionsverweis {medan-medan normal satuan}{}{}
dalam kasus ini.

}
{} {}

\inputaufgabegibtloesung
{}
{

Misalkan
\mathkor {} {r} {dan} {R} {}
bilangan real positif yang memenuhi

\mavergleichskette
{\vergleichskette
{ 0
}
{ < }{ r
}
{ < }{ R
}
{    }{
}
{   }{
}
}
{}{}{.}
Tentukan suatu
\definitionsverweis {medan normal satuan}{}{}
untuk torus
\zusatzklammer {tertanam} {} {}

\mavergleichskettedisp
{\vergleichskette
{ T
}
{ =} { \left\{ (x,y,z) \in \R^3  \mid   \left(  \sqrt{x^2+y^2} -R  \right)^2 +z^2  = r^2         \right\}
}
{ } {
}
{ } {
}
{ } {
}
}
{}{}{.}

}
{} {}

\inputaufgabe
{}
{

Misalkan
\maabbdisp {f} { I } {\R_+
} {}
sebuah
\definitionsverweis {fungsi terdiferensialkan}{}{}
dan misalkan $Y$ adalah permukaan putar yang terkait dengan grafik $f$, dipandang sebagai hipermuka terdiferensialkan di $\R^3$ dalam pengertian
Lemma 2.1.
Tentukan
\definitionsverweis {medan-medan normal satuan}{}{}
dalam kasus ini.

}
{} {}

\inputaufgabe
{}
{

Buktikan bahwa
\definitionsverweis {pemetaan Gauss}{}{}
pada sebuah hiperbidang bersifat konstan. Apakah pernyataan sebaliknya juga berlaku?

}
{} {}

\inputaufgabe
{}
{

Buktikan bahwa
\definitionsverweis {pemetaan Gauss}{}{}
pada sebuah permukaan bola di $\R^n$ yang berpusat di titik asal, untuk setiap pilihan orientasi, diinduksi oleh homoteti skalar pada $\R^n$ (dengan faktor positif untuk orientasi keluar dan faktor negatif untuk orientasi masuk).

}
{} {}

\inputaufgabe
{}
{

Misalkan

\mavergleichskette
{\vergleichskette
{ Y
}
{ \subseteq }{ \R^n
}
{  }{
}
{    }{
}
{   }{
}
}
{}{}{}
sebuah
\definitionsverweis {hipermuka terdiferensialkan}{}{}
dengan
\definitionsverweis {medan normal satuan}{}{}
$N$ dan medan normal satuan yang berlawanan, $-N$. Buktikan bahwa kedua
\definitionsverweis {pemetaan Gauss}{}{}
yang terkait saling berhubungan melalui
\definitionsverweis {pemetaan antipodal}{}{}
\maabbeledisp {} {S^{n-1}} {S^{n-1}
} { P } { -P
} {.}

}
{} {}

\inputaufgabegibtloesung
{}
{

Misalkan

\mavergleichskette
{\vergleichskette
{ \alpha,\beta
}
{ \in }{ \R_+
}
{  }{
}
{    }{
}
{   }{
}
}
{}{}{}
dan misalkan

\mavergleichskettedisp
{\vergleichskette
{ Y
}
{ =} { \left\{ (x,y)\in\R^2  \mid   \alpha x^2+\beta y^2 = 1        \right\}
}
{ } {
}
{ } {
}
{ } {
}
}
{}{}{.}
Orientasikan $Y$ dengan medan normal satuan yang diperoleh dari gradien keluar fungsi
$ (x,y) \longmapsto \alpha x^2+\beta y^2 $.
\aufzaehlungzwei {Tentukan
\definitionsverweis {pemetaan Gauss}{}{}
untuk $Y$.
} { Berikan secara eksplisit suatu
\definitionsverweis {pemetaan invers}{}{}
untuk pemetaan Gauss pada $Y$.
}

}
{} {}

\inputaufgabegibtloesung
{}
{

Misalkan

\mavergleichskette
{\vergleichskette
{ \alpha,\beta
}
{ \in }{ \R_+
}
{  }{
}
{    }{
}
{   }{
}
}
{}{}{}
dan misalkan

\mavergleichskettedisp
{\vergleichskette
{ Y
}
{ =} { \left\{ (x,y)\in\R^2  \mid   \alpha x^2+\beta y^2 = 1        \right\}
}
{ } {
}
{ } {
}
{ } {
}
}
{}{}{.}
Buktikan bahwa
\definitionsverweis {pemetaan Gauss}{}{}
untuk $Y$ bersifat bijektif.

}
{} {}

\inputaufgabe
{}
{

Berikan sebuah contoh
\definitionsverweis {hipermuka terdiferensialkan}{}{}

\mavergleichskette
{\vergleichskette
{ Y
}
{ \subseteq }{ \R^2
}
{  }{
}
{    }{
}
{   }{
}
}
{}{}{}
sedemikian sehingga
\definitionsverweis {pemetaan Gauss}{}{}
bersifat surjektif dan terdapat titik-titik pada $S^1$ yang dicapai tak berhingga kali.

}
{} {}

\inputaufgabe
{}
{

Misalkan

\mavergleichskette
{\vergleichskette
{ W
}
{ \subseteq }{ \R^n
}
{  }{
}
{    }{
}
{   }{
}
}
{}{}{}
\definitionsverweis {terbuka}{}{,}
\maabb {h} { W } { \R
} {}
sebuah
\definitionsverweis {fungsi yang terdiferensialkan secara kontinu}{}{}
dan

\mavergleichskette
{\vergleichskette
{ Y
}
{ = }{ h^{-1}(c)
}
{  }{
}
{    }{
}
{   }{
}
}
{}{}{}
adalah
\definitionsverweis {serat}{}{}
di atas

\mavergleichskette
{\vergleichskette
{ c
}
{ \in }{ \R
}
{  }{
}
{    }{
}
{   }{
}
}
{}{}{,}
dengan $h$
\definitionsverweis {reguler}{}{}
di setiap titik pada $Y$. Misalkan
\maabbdisp {L} {\R^n } { \R^n
} {}
sebuah
\definitionsverweis {isometri linear}{}{,}

\mavergleichskette
{\vergleichskette
{ W'
}
{ = }{ L^{-1}(W)
}
{  }{
}
{    }{
}
{   }{
}
}
{}{}{}
dan

\mavergleichskettedisp
{\vergleichskette
{ Z
}
{ =} { L^{-1}(Y)
}
{ =} { (h \circ L)^{-1}(c)
}
{ } {
}
{ } {
}
}
{}{}{.}
Buktikan bahwa konsep-konsep kurva geodesik, medan normal, medan normal satuan, dan pemetaan Gauss pada $Y$ dan $Z$ saling bersesuaian
\zusatzklammer {yakni dapat dipetakan satu sama lain dengan bantuan $L$} {} {.}

}
{} {}

  \zwischenueberschrift{Soal untuk dikumpulkan}

\inputaufgabe
{3}
{

Tentukan citra parabola standar berorientasi ke atas di bawah
\definitionsverweis {pemetaan Gauss}{}{.}

}
{} {}

\inputaufgabe
{4}
{

Misalkan

\mavergleichskette
{\vergleichskette
{ W
}
{ = }{ \R^3 \setminus \{ (0,0,0) \}
}
{  }{
}
{    }{
}
{   }{
}
}
{}{}{}
dan misalkan

\mavergleichskette
{\vergleichskette
{ Y
}
{ \subseteq }{ W
}
{  }{
}
{    }{
}
{   }{
}
}
{}{}{}
adalah
\definitionsverweis {himpunan nol}{}{}
dari

\mavergleichskettedisp
{\vergleichskette
{ h(x,y,z)
}
{ =} { x^2+y^2-z^2
}
{ } {
}
{ } {
}
{ } {
}
}
{}{}{.}
Tentukan citra $Y$
\zusatzklammer {dengan orientasi yang ditentukan oleh gradien $h$} {} {}
di bawah
\definitionsverweis {pemetaan Gauss}{}{.}

}
{} {}

\inputaufgabe
{5}
{

Misalkan

\mavergleichskette
{\vergleichskette
{ \alpha,\beta, \gamma
}
{ \in }{ \R_+
}
{  }{
}
{    }{
}
{   }{
}
}
{}{}{}
dan misalkan

\mavergleichskettedisp
{\vergleichskette
{ Y
}
{ =} { \left\{ (x,y,z) \in \R^3  \mid   \alpha x^2+ \beta y^2 + \gamma z^2 = 1        \right\}
}
{ } {
}
{ } {
}
{ } {
}
}
{}{}{.}
Buktikan bahwa
\definitionsverweis {pemetaan Gauss}{}{}
untuk $Y$ bersifat bijektif.

}
{} {}

\inputaufgabe
{3}
{

Misalkan

\mavergleichskette
{\vergleichskette
{ W
}
{ \subseteq }{ \R^n
}
{  }{
}
{    }{
}
{   }{
}
}
{}{}{}
\definitionsverweis {terbuka}{}{,}
\maabb {h} { W } { \R
} {}
sebuah
\definitionsverweis {fungsi yang terdiferensialkan secara kontinu}{}{}
dan

\mavergleichskette
{\vergleichskette
{ Y
}
{ = }{ h^{-1}(c)
}
{  }{
}
{    }{
}
{   }{
}
}
{}{}{}
adalah
\definitionsverweis {serat}{}{}
di atas

\mavergleichskette
{\vergleichskette
{ c
}
{ \in }{ \R
}
{  }{
}
{    }{
}
{   }{
}
}
{}{}{,}
dengan $h$
\definitionsverweis {reguler}{}{}
di setiap titik pada $Y$. Misalkan
\maabbdisp {L} {\R^n } { \R^n
} {}
sebuah
\definitionsverweis {isomorfisme linear}{}{,}

\mavergleichskette
{\vergleichskette
{ W'
}
{ = }{ L^{-1}(W)
}
{  }{
}
{    }{
}
{   }{
}
}
{}{}{}
dan

\mavergleichskettedisp
{\vergleichskette
{ Z
}
{ =} { L^{-1}(Y)
}
{ =} { (h \circ L)^{-1}(c)
}
{ } {
}
{ } {
}
}
{}{}{.}
Buktikan bahwa konsep-konsep kurva geodesik, medan normal, medan normal satuan, dan pemetaan Gauss pada $Y$ dan $Z$ secara umum tidak saling bersesuaian
\zusatzklammer {berbeda dengan
Soal 2.14,
yang melibatkan suatu isometri} {} {.}

}
{} {}
